\documentclass[10pt,a4paper]{amsart}
\usepackage{amsmath,amssymb,amsthm,graphicx}
\usepackage{verbatim}
\newenvironment{reference}{\comment}{\endcomment}
%\newenvironment{reference}{}{}
\newenvironment{slogan}{\comment}{\endcomment}
\newenvironment{history}{\comment}{\endcomment}

% For commutative diagrams we use Xy-pic
\usepackage[all]{xy}

% We use 2cell for 2-commutative diagrams.
\xyoption{2cell}
% Theorem environments.
%
\theoremstyle{plain}
\newtheorem{theorem}[subsection]{Theorem}
\newtheorem{proposition}[subsection]{Proposition}
\newtheorem{lemma}[subsection]{Lemma}

\theoremstyle{definition}
\newtheorem{definition}[subsection]{Definition}
\newtheorem{example}[subsection]{Example}
\newtheorem{exercise}[subsection]{Exercise}
\newtheorem{situation}[subsection]{Situation}

\theoremstyle{remark}
\newtheorem{remark}[subsection]{Remark}
\newtheorem{remarks}[subsection]{Remarks}

\numberwithin{equation}{subsection}

% Macros
%
\def\lim{\mathop{\mathrm{lim}}\nolimits}
\def\colim{\mathop{\mathrm{colim}}\nolimits}
\def\Spec{\mathop{\mathrm{Spec}}}
\def\Hom{\mathop{\mathrm{Hom}}\nolimits}
\def\Ext{\mathop{\mathrm{Ext}}\nolimits}
\def\SheafHom{\mathop{\mathcal{H}\!\mathit{om}}\nolimits}
\def\SheafExt{\mathop{\mathcal{E}\!\mathit{xt}}\nolimits}
\def\Sch{\mathit{Sch}}
\def\Mor{\mathop{\mathrm{Mor}}\nolimits}
\def\Ob{\mathop{\mathrm{Ob}}\nolimits}
\def\Sh{\mathop{\mathit{Sh}}\nolimits}
\def\NL{\mathop{N\!L}\nolimits}
\def\CH{\mathop{\mathrm{CH}}\nolimits}
\def\proetale{{pro\text{-}\acute{e}tale}}
\def\etale{{\acute{e}tale}}
\def\QCoh{\mathit{QCoh}}
\def\Ker{\mathop{\mathrm{Ker}}}
\def\Im{\mathop{\mathrm{Im}}}
\def\Coker{\mathop{\mathrm{Coker}}}
\def\Coim{\mathop{\mathrm{Coim}}}

% Boxtimes
%
\DeclareMathSymbol{\boxtimes}{\mathbin}{AMSa}{"02}

%
% Macros for moduli stacks/spaces
%
\def\QCohstack{\mathcal{QC}\!\mathit{oh}}
\def\Cohstack{\mathcal{C}\!\mathit{oh}}
\def\Spacesstack{\mathcal{S}\!\mathit{paces}}
\def\Quotfunctor{\mathrm{Quot}}
\def\Hilbfunctor{\mathrm{Hilb}}
\def\Curvesstack{\mathcal{C}\!\mathit{urves}}
\def\Polarizedstack{\mathcal{P}\!\mathit{olarized}}
\def\Complexesstack{\mathcal{C}\!\mathit{omplexes}}
% \Pic is the operator that assigns to X its picard group, usage \Pic(X)
% \Picardstack_{X/B} denotes the Picard stack of X over B
% \Picardfunctor_{X/B} denotes the Picard functor of X over B
\def\Pic{\mathop{\mathrm{Pic}}\nolimits}
\def\Picardstack{\mathcal{P}\!\mathit{ic}}
\def\Picardfunctor{\mathrm{Pic}}
\def\Deformationcategory{\mathcal{D}\!\mathit{ef}}
\usepackage{cleveref}
\crefname{originallemma}{original lemma}{original lemmas}
\Crefname{originallemma}{Original lemma}{Original lemmas}
\makeatletter
\newlabel{algebra-lemma-pullback-module}{{02M1}{source}{}{sourcepointer-1}{}}
\newlabel{algebra-lemma-pullback-module@cref}{{[originallemma][0][]02M1}{[1][0][]source}{}{}{}}
\makeatother
\crefname{originallemma}{original lemma}{original lemmas}
\Crefname{originallemma}{Original lemma}{Original lemmas}
\makeatletter
\newlabel{chow-lemma-additivity-divisors-restricted}{{02QG}{source}{}{sourcepointer-2}{}}
\newlabel{chow-lemma-additivity-divisors-restricted@cref}{{[originallemma][0][]02QG}{[1][0][]source}{}{}{}}
\makeatother
\crefname{originalremark}{original remark}{original remarks}
\Crefname{originalremark}{Original remark}{Original remarks}
\makeatletter
\newlabel{decent-spaces-remark-functoriality-henselian-local-ring}{{0EPL}{source}{}{sourcepointer-3}{}}
\newlabel{decent-spaces-remark-functoriality-henselian-local-ring@cref}{{[originalremark][0][]0EPL}{[1][0][]source}{}{}{}}
\makeatother
\crefname{originalreference}{original reference}{original references}
\Crefname{originalreference}{Original reference}{Original references}
\makeatletter
\newlabel{definition-flat-pullback}{{02RB}{source}{}{sourcepointer-4}{}}
\newlabel{definition-flat-pullback@cref}{{[originalreference][0][]02RB}{[1][0][]source}{}{}{}}
\makeatother
\crefname{originalreference}{original reference}{original references}
\Crefname{originalreference}{Original reference}{Original references}
\makeatletter
\newlabel{lemma-etale-pullback}{{0EPY}{source}{}{sourcepointer-5}{}}
\newlabel{lemma-etale-pullback@cref}{{[originalreference][0][]0EPY}{[1][0][]source}{}{}{}}
\makeatother
\crefname{originalreference}{original reference}{original references}
\Crefname{originalreference}{Original reference}{Original references}
\makeatletter
\newlabel{lemma-proper-image}{{0EEG}{source}{}{sourcepointer-6}{}}
\newlabel{lemma-proper-image@cref}{{[originalreference][0][]0EEG}{[1][0][]source}{}{}{}}
\makeatother
\crefname{originalreference}{original reference}{original references}
\Crefname{originalreference}{Original reference}{Original references}
\makeatletter
\newlabel{remark-irreducible-component}{{0EE8}{source}{}{sourcepointer-7}{}}
\newlabel{remark-irreducible-component@cref}{{[originalreference][0][]0EE8}{[1][0][]source}{}{}{}}
\makeatother
\crefname{originalreference}{original reference}{original references}
\Crefname{originalreference}{Original reference}{Original references}
\makeatletter
\newlabel{situation-setup}{{02QL}{source}{}{sourcepointer-8}{}}
\newlabel{situation-setup@cref}{{[originalreference][0][]02QL}{[1][0][]source}{}{}{}}
\makeatother
\crefname{originalsection}{original section}{original sections}
\Crefname{originalsection}{Original section}{Original sections}
\makeatletter
\newlabel{spaces-morphisms-section-rational-maps}{{0EML}{source}{}{sourcepointer-9}{}}
\newlabel{spaces-morphisms-section-rational-maps@cref}{{[originalsection][0][]0EML}{[1][0][]source}{}{}{}}
\makeatother
\title{Pullback of a principal divisor along an etale morphism is the sum over the irreducible components of the pushforwards of the restricted divisors}
\author{The Stacks Project Authors}
\date{stacks-project:spaces-chow.tex:local-snapshot-20261012; GNU FDL 1.2 or later}
\begin{document}
\maketitle
\noindent\textit{Editorial note.} All delivered mathematical text is the author's own, copied from the pinned Stacks Project chapter file spaces-chow.tex; no statement or proof step was rewritten, reordered or removed.\par\medskip
\noindent\textit{Editorial note.} The counted proof is the author's complete argument: a cycle-theoretic reduction to a coefficient comparison at generic points in henselizations, followed by a length computation over total rings of fractions that is finished by an additivity statement and a disjoint-union decomposition\par\medskip
\noindent\textit{Editorial note.} The label lemma-etale-pullback-principal-divisor names the counted statement, and the proof environment immediately following it in the source is delivered verbatim as its proof.\par\medskip
\noindent\textit{Editorial note.} The counted statement is carried by the author's own lemma environment in the chapter source.\par\medskip
\noindent\textit{Allowed standard prerequisites (editorial).}\begin{itemize}
\item Flat pullback of cycles with relative dimension 0, and the definition of the k-dimensional part of a cycle, so that the pullback of a divisor by an etale morphism is defined
\item Deltas of points and their behaviour for etale morphisms, the generic point of an integral closed subspace and the scheme theoretic image construction for the induced closed subspace of the base
\item Henselizations of reduced Noetherian local rings, their total rings of fractions, and the length formula for the order of a rational function along an integral closed subspace of codimension 1
\end{itemize}\medskip
\ifdefined\hypertarget\hypertarget{sourcepointer-1}{}\fi
\noindent\textit{Source reference.} Original lemma 02M1: sources/stacks-chow/tags.txt, source lines 3428--3428. Prerequisite pointer invoked at proof line(s) 1777; it names a standard fact of a sibling chapter that the counted argument uses, and every step of the counted proof built on it is retained verbatim in the delivered body.\par\medskip
\ifdefined\hypertarget\hypertarget{sourcepointer-2}{}\fi
\noindent\textit{Source reference.} Original lemma 02QG: sources/stacks-chow/tags.txt, source lines 3587--3587. Prerequisite pointer invoked at proof line(s) 1787; it names a standard fact of a sibling chapter that the counted argument uses, and every step of the counted proof built on it is retained verbatim in the delivered body.\par\medskip
\ifdefined\hypertarget\hypertarget{sourcepointer-3}{}\fi
\noindent\textit{Source reference.} Original remark 0EPL: sources/stacks-chow/tags.txt, source lines 18381--18381. Prerequisite pointer invoked at proof line(s) 1726; it names a standard fact of a sibling chapter that the counted argument uses, and every step of the counted proof built on it is retained verbatim in the delivered body.\par\medskip
\ifdefined\hypertarget\hypertarget{sourcepointer-4}{}\fi
\noindent\textit{Source reference.} Original reference 02RB: sources/stacks-spaces-chow/pinned-original-source.tex, source lines 1262--1262. Prerequisite pointer invoked at proof line(s) 1701; it names a standard fact of the same chapter that the counted argument uses, and every step of the counted proof built on it is retained verbatim in the delivered body.\par\medskip
\ifdefined\hypertarget\hypertarget{sourcepointer-5}{}\fi
\noindent\textit{Source reference.} Original reference 0EPY: sources/stacks-spaces-chow/pinned-original-source.tex, source lines 1332--1332. Prerequisite pointer invoked at proof line(s) 1742; it names a standard fact of the same chapter that the counted argument uses, and every step of the counted proof built on it is retained verbatim in the delivered body.\par\medskip
\ifdefined\hypertarget\hypertarget{sourcepointer-6}{}\fi
\noindent\textit{Source reference.} Original reference 0EEG: sources/stacks-spaces-chow/pinned-original-source.tex, source lines 820--820. Prerequisite pointer invoked at proof line(s) 1715; it names a standard fact of the same chapter that the counted argument uses, and every step of the counted proof built on it is retained verbatim in the delivered body.\par\medskip
\ifdefined\hypertarget\hypertarget{sourcepointer-7}{}\fi
\noindent\textit{Source reference.} Original reference 0EE8: sources/stacks-spaces-chow/pinned-original-source.tex, source lines 650--650. Prerequisite pointer invoked at proof line(s) 1689; it names a standard fact of the same chapter that the counted argument uses, and every step of the counted proof built on it is retained verbatim in the delivered body.\par\medskip
\ifdefined\hypertarget\hypertarget{sourcepointer-8}{}\fi
\noindent\textit{Source reference.} Original reference 02QL: sources/stacks-spaces-chow/pinned-original-source.tex, source lines 264--264. Prerequisite pointer invoked at proof line(s) 1681; it names a standard fact of the same chapter that the counted argument uses, and every step of the counted proof built on it is retained verbatim in the delivered body.\par\medskip
\ifdefined\hypertarget\hypertarget{sourcepointer-9}{}\fi
\noindent\textit{Source reference.} Original section 0EML: sources/stacks-chow/tags.txt, source lines 18311--18311. Prerequisite pointer invoked at proof line(s) 1697; it names a standard fact of a sibling chapter that the counted argument uses, and every step of the counted proof built on it is retained verbatim in the delivered body.\par\medskip
\begin{lemma}
\label{lemma-etale-pullback-principal-divisor}
In Situation \ref{situation-setup} let $f : X \to Y$ be an \'etale
morphism of good algebraic spaces over $B$. Assume $Y$ is integral.
Let $g \in R(Y)^*$. As cycles on $X$ we have
$$
f^*(\text{div}_Y(g)) =
\sum\nolimits_{X'} (X' \to X)_*\text{div}_{X'}(g \circ f|_{X'})
$$
where the sum is over the irreducible components of $X$
(Remark \ref{remark-irreducible-component}).
\end{lemma}

\begin{proof}
The map $|X| \to |Y|$ is open. The set of irreducible components of $|X|$
is locally finite in $|X|$. We conclude that $f|_{X'} : X' \to Y$
is dominant for every irreducible component $X' \subset X$.
Thus $g \circ f|_{X'}$ is defined (Morphisms of Spaces, Section
\ref{spaces-morphisms-section-rational-maps}), hence
$\text{div}_{X'}(g \circ f|_{X'})$ is defined. Moreover, the
sum is locally finite and we find that the right hand side indeed
is a cycle on $X$. The left hand side is defined by
Definition \ref{definition-flat-pullback}
and the fact that an \'etale morphism is flat of relative dimension $0$.

\medskip\noindent
Since $f$ is \'etale we see that $\delta_X(x) = \delta_y(f(x))$
for all $x \in |X|$. Thus if $\dim_\delta(Y) = n$, then $\dim_\delta(X') = n$
for every irreducible component $X'$ of $X$ (since generic points of $X$
map to the generic point of $Y$, see above). Thus both left
and right hand side are $(n - 1)$-cycles.

\medskip\noindent
Let $Z \subset X$ be an integral closed subspace with $\dim_\delta(Z) = n - 1$.
To prove the equality, we need to show that the coefficients
of $Z$ are the same. Let $Z' \subset Y$ be the integral closed
subspace constructed in Lemma \ref{lemma-proper-image}.
Then $\dim_\delta(Z') = n - 1$ too. Let $\xi \in |Z|$ be the generic point.
Then $\xi' = f(\xi) \in |Z'|$ is the generic point.
Consider the commutative diagram
$$
\xymatrix{
\Spec(\mathcal{O}_{X, \xi}^h) \ar[r] \ar[d] & X \ar[d] \\
\Spec(\mathcal{O}_{Y, \xi'}^h) \ar[r] & Y
}
$$
of Decent Spaces, Remark
\ref{decent-spaces-remark-functoriality-henselian-local-ring}.
We have to be slightly careful as the reduced Noetherian local rings
$\mathcal{O}_{X, \xi}^h$ and $\mathcal{O}_{Y, \xi'}^h$ need not be domains.
Thus we work with total rings of fractions $Q(-)$ rather than fraction fields.
By definition, to get the coefficient of $Z'$ in $\text{div}_Y(g)$
we write the image of $g$ in $Q(\mathcal{O}_{Y, \xi'}^h)$
as $a/b$ with $a, b \in \mathcal{O}_{Y, \xi'}^h$ nonzerodivisors
and we take
$$
\text{ord}_{Z'}(g) =
\text{length}_{\mathcal{O}_{Y, \xi'}^h}
(\mathcal{O}_{Y, \xi'}^h/a \mathcal{O}_{Y, \xi'}^h) -
\text{length}_{\mathcal{O}_{Y, \xi'}^h}
(\mathcal{O}_{Y, \xi'}^h/b \mathcal{O}_{Y, \xi'}^h)
$$
Observe that the coefficient of $Z$ in $f^*\text{div}_Y(G)$
is the same integer, see Lemma \ref{lemma-etale-pullback}.
Suppose that $\xi \in X'$. Then we can consider the maps
$$
\mathcal{O}_{Y, \xi'}^h \to
\mathcal{O}_{X, \xi}^h \to
\mathcal{O}_{X', \xi}^h
$$
The first arrow is flat and the second arrow is a surjective
map of reduced local Noetherian rings of dimension $1$.
Therefore both these maps send nonzerodivisors to nonzerodivisors
and we conclude the coefficient of $Z'$ in $\text{div}_{X'}(g \circ f|_{X'})$
is
$$
\text{ord}_Z(g \circ f|_{X'}) =
\text{length}_{\mathcal{O}_{X', \xi}^h}
(\mathcal{O}_{X', \xi}^h/a \mathcal{O}_{X', \xi}^h) -
\text{length}_{\mathcal{O}_{Y, \xi'}^h}
(\mathcal{O}_{X', \xi}^h/b \mathcal{O}_{X', \xi}^h)
$$
by the same prescription as above. Thus it suffices to show
$$
\text{length}_{\mathcal{O}_{Y, \xi'}^h}
(\mathcal{O}_{Y, \xi'}^h/a \mathcal{O}_{Y, \xi'}^h) =
\sum\nolimits_{\xi \in |X'|}
\text{length}_{\mathcal{O}_{X', \xi}^h}
(\mathcal{O}_{X', \xi}^h/a \mathcal{O}_{X', \xi}^h)
$$
First, since the ring map $\mathcal{O}_{Y, \xi'}^h \to
\mathcal{O}_{X, \xi}^h$ is flat and unramified, we have
$$
\text{length}_{\mathcal{O}_{Y, \xi'}^h}
(\mathcal{O}_{Y, \xi'}^h/a \mathcal{O}_{Y, \xi'}^h) =
\text{length}_{\mathcal{O}_{X, \xi}^h}
(\mathcal{O}_{X, \xi}^h/a \mathcal{O}_{X, \xi}^h)
$$
by Algebra, Lemma \ref{algebra-lemma-pullback-module}.
Let $\mathfrak q_1, \ldots, \mathfrak q_t$ be the nonmaximal
primes of $\mathcal{O}_{X, \xi}^h$ and set
$R_j = \mathcal{O}_{X, \xi}^h/\mathfrak q_j$.
For $X'$ as above, denote $J(X') \subset \{1, \ldots, t\}$
the set of indices such that $\mathfrak q_j$ corresponds
to a point of $X'$, i.e., such that under the surjection
$\mathcal{O}_{X, \xi}^h \to \mathcal{O}_{X', \xi}$
the prime $\mathfrak q_j$ corresponds to a prime
of $\mathcal{O}_{X', \xi}$.
By Chow Homology, Lemma \ref{chow-lemma-additivity-divisors-restricted}
we get
$$
\text{length}_{\mathcal{O}_{X, \xi}^h}
(\mathcal{O}_{X, \xi}^h/a \mathcal{O}_{X, \xi}^h) =
\sum\nolimits_j \text{length}_{R_j}(R_j/a R_j)
$$
and
$$
\text{length}_{\mathcal{O}_{X', \xi}^h}
(\mathcal{O}_{X', \xi}^h/a \mathcal{O}_{X', \xi}^h) =
\sum\nolimits_{j \in J(X')} \text{length}_{R_j}(R_j/a R_j)
$$
Thus the result of the lemma holds because
$\{1, \ldots, t\}$ is the disjoint union of the sets $J(X')$:
each point of codimension $0$ on $X$ lies on a unique $X'$.
\end{proof}

\end{document}
